% ---------------------------------------------------------------------------
%  Formal table of where the reference open-set implementation leaves the
%  test-time setting. Same format as the paper/code divergence table.
%  Source: github.com/ispc-lab/LEAD, inspected at HEAD.
% ---------------------------------------------------------------------------
\subsection{Where the reference implementation leaves the test-time setting}
\label{sec:os-lead}

The open-set and universal domain adaptation methods we build on are distributed with
implementations that assume offline access to the target domain. That assumption is not
stated as a limitation, and it is not visible from the papers, but it decides what the
reported numbers mean. Table~\ref{tab:leadviol} records the four places where the
reference implementation of LEAD \cite{lead} departs from the Tier-1 definition of
Section~\ref{sec:method-tiers}, with the line that does it.

\begin{table*}[t]
\centering
\caption{The reference LEAD implementation (\texttt{ispc-lab/LEAD}) against the Tier-1
test-time constraints. Line numbers are from \texttt{train\_target.py} unless marked.
Together these make the published protocol source-free domain adaptation rather than
test-time adaptation: it requires the whole target set in advance, revisits it fifty
times, and evaluates on the data it adapted on.}
\label{tab:leadviol}
\small
\begin{tabular}{@{}p{0.20\textwidth}p{0.17\textwidth}p{0.26\textwidth}p{0.29\textwidth}@{}}
\toprule
Mechanism & Code & Constraint violated & Consequence \\
\midrule
Multi-epoch offline training on the target set
 & \texttt{epochs=50} (\texttt{config/model\_config.py}:15); loop at :375
 & Adaptation must be a single online pass over the stream
 & The whole target set must exist before adaptation starts, and every sample is revisited
   fifty times. No deployed system has this. \\
\addlinespace[3pt]
Adaptation and evaluation share one dataset
 & :360--362, both \texttt{DataLoader}s wrap the same \texttt{target\_dataset}
 & Evaluation must not be adapted on
 & Every sample scored at :383 was trained on at :377. The reported score is training
   performance, and no held-out target split exists anywhere in the file. \\
\addlinespace[3pt]
Pseudo-labels built from a full pass over the evaluation set
 & :209 calls \texttt{obtain\_LEAD\_pseudo\_labels} on \texttt{test\_dataloader} at the
   start of every epoch
 & Predictions must depend only on samples already seen
 & The model sees the entire evaluation set, including its own future, before adapting.
   Known- and unknown-space bases are fitted on data not yet encountered. \\
\addlinespace[3pt]
Stream order discarded
 & \texttt{shuffle=True} on the adaptation loader, :360
 & The target stream must be processed in acquisition order
 & Temporal structure is destroyed and each batch mixes the whole recording period, so
   batch statistics are those of the full target distribution rather than of the moment. \\
\bottomrule
\end{tabular}
\end{table*}

None of the four is a defect of the method as described; each is a property of the
evaluation harness it ships with. The distinction matters because the first and third in
particular make the reported figures an upper bound that a streaming deployment cannot
approach, and because a reimplementation that honoured the constraints would not be
comparable with the published numbers. All open-set results in
Section~\ref{sec:os-methods} are therefore reported against our own Tier-1 source-only
baseline rather than against published figures obtained under this harness.
